\chapter{Principe Extrémal}

	\section{Quelques principes}
Le principe extrémal consiste à considérer, parmi les objets étudiés, celui qui est le plus petit ou le plus grand pour un critère choisi (longueur, aire, valeur\ldots), puis à exploiter cette propriété extrême. Encore faut-il que cet objet existe : c'est ce que garantissent les faits suivants.
\begin{enumerate}
	\item Tout ensemble fini non vide de nombres réels admet un plus petit élément et un plus grand élément. 
	\item Tout sous-ensemble non vide de $\mathbb{N}$ admet un plus petit élément. (C'est faux pour un ensemble de réels en général : $]0\,;1]$ n'a pas de plus petit élément.)
	\item Soit $A$ un ensemble non vide de nombres réels (éventuellement infini). 
	\begin{itemize}
		\item On dit que $A$ est \emph{majoré} s'il existe un réel $M$ tel que $a \leq M$ pour tout $a \in A$ ; un tel $M$ est un \emph{majorant} de $A$. Si $A$ est majoré, alors il a un plus petit majorant, mais pas nécessairement un plus grand élément (exemple : $[0\,;1[$ a pour plus petit majorant $1$, qui n'est pas dans l'ensemble). 
		\item On dit que $A$ est \emph{minoré} s'il existe un réel $m$ tel que $m \leq a$ pour tout $a \in A$ ; un tel $m$ est un \emph{minorant} de $A$. Si $A$ est minoré, alors il a un plus grand minorant, mais pas nécessairement un plus petit élément (exemple : $]0\,;1]$ a pour plus grand minorant $0$, qui n'est pas dans l'ensemble).
	\end{itemize}
\end{enumerate}
La \emph{descente infinie} est une conséquence du deuxième fait : une suite strictement décroissante d'entiers naturels ne peut pas être infinie. En effet, si elle l'était, l'ensemble de ses termes serait une partie non vide de $\mathbb{N}$ sans plus petit élément, car chaque terme est suivi d'un terme strictement plus petit.


\section{Exercices}
\begin{exercise}
	\begin{enumerate}
		\item Prouver qu'il n'existe pas de couple d'entiers strictement positifs $(x,y)$ tels que : $x^2 = 2 y^2$.
		\item En déduire que $\sqrt{2}$ est irrationnel.
	\end{enumerate}
\end{exercise}

\begin{exercise}[E9]
	Prouver qu'il n'y a pas de quadruplets d'entiers strictement positifs $(x,y,a,b)$ tels que :
	\[
	x^2 + y^2 = 3 (a^2 + b^2).
	\]
\end{exercise}

\begin{exercise}[E3]
	Considérons $n \geq 3$ points distincts du plan tels que trois quelconques d'entre eux forment un triangle d'aire $\leq 1$ (un triangle aplati a une aire nulle).
	Montrer que les $n$ points se trouvent dans un triangle d'aire $\leq 4$.
\end{exercise}

\begin{exercise}[E10]
	Considérons $n$ points distincts du plan qui ne sont pas tous sur une même droite. 
	Montrez qu'il existe une droite qui passe par exactement deux de ces points (et par aucun autre). 
\end{exercise}

\begin{exercise}[E5]
	Soit $P$ un ensemble non vide de points du plan. 
	Tout point de $P$ est le milieu de deux points \emph{distincts} de $P$.
	Montrer que $P$ est un ensemble infini.
\end{exercise}

\section{Problèmes}
\begin{problem}[12]
	Il y a 20 pays sur une planète.
	Deux pays ont une relation diplomatique si et seulement si chacun possède une ambassade dans l'autre ; il n'y a pas d'autre ambassade.
	Pour chaque groupe de trois pays, il y en a deux qui n'ont aucune relation diplomatique.
	Montrer qu'il y a au plus 200 ambassades sur cette planète.  
\end{problem}

\begin{problem}[23]
	Prenez trois points $A_1$, $B_1$, $C_1$, respectivement sur les côtés $BC$, $CA$, $AB$ d'un triangle $ABC$.
	Montrer que si $AA_1 \leq 1$, $BB_1 \leq 1$, $CC_1 \leq 1$, alors l'aire du triangle est $\leq \frac{1}{\sqrt{3}}$.
\end{problem}

\begin{problem}[21]
	Montrer que tout polygone convexe d'aire 1 est contenu dans un rectangle d'aire 2. 
\end{problem}

\begin{problem}[28]
	On choisit 8 points dans un disque (bord compris) de rayon 1.
	Montrer qu'il y a deux points distants d'au plus 1.
\end{problem}

\begin{problem}[32]
	Déterminer toutes les solutions en réels strictement positifs du système d'équations :
	\[
	x_1 + x_2 = x_3^2,
	\quad
	x_2 + x_3 = x_4^2,
	\quad
	x_3 + x_4 = x_5^2,
	\quad
	x_4 + x_5 = x_1^2,
	\quad
	x_5 + x_1 = x_2^2.
	\]
\end{problem}